% Autor: René Henke, David Steffen, Miriam Neumann
% Datum: 2015-08-19

\section[Die junge DPG stellt sich vor]{Mehr als nur Uni -- Die junge~DPG stellt sich vor}
\begin{minipage}{0.1\textwidth}
\includegraphics[width=\textwidth]{res/jdpg_logo.pdf}
\end{minipage}\hfill
\begin{minipage}{0.75\textwidth}
\centering
\textbf{Liebe Physik-Erstis der Uni Münster,\\
auch die junge DPG möchte Euch herzlich im Kreis der Physik-Studenten begrüßen.}
\end{minipage}\hfill
\begin{minipage}{0.1\textwidth}
\includegraphics[width=\textwidth]{res/jdpg_logo.pdf}
\end{minipage}

\begin{multicols}{2}
In Deutschland studieren rund \num{30000} Studenten Physik. Du bist jetzt einer von ihnen. An "Jungphysiker" wie dich wendet sich die junge~Deutsche~Physikalische~Gesellschaft~(jDPG) mit ihren Angeboten rund um das Physikstudium.

Die Ziele der jDPG sind unter anderem ein Netzwerk für Studierende in ganz Deutschland zu schaffen, sowie ein wissenschaftliches Programm zu bieten. Desweiteren ermöglicht die jDPG eine umfangreiche Berufsvorbereitung und fördert den Dialog zwischen den Generationen. Dabei versteht sie sich, zusammen mit den Physikfachschaften, als eine bundesweite Interessenvertretung von Physikstudenten.

Mehr als \num{3000} Mitglieder zählt die jDPG derzeit. Für diese Mitglieder und auch für alle anderen Physik-Studierenden gibt es unsere Angebote sowohl deutschlandweit, als auch vor Ort in den einzelnen Regionalgruppen.

Die Regionalgruppe Münster veranstaltet jedes Semester eine Exkursion zu Physikalisch interessanten Zielen (Forschungsinstitute, Unternehmen, ect.). So ging es in den letzten Jahren z.\,B.\ in das Endlager KONRAD, zum Kernkraftwerk in Lingen, sowie zum weltweit ersten Supraleiterkabel, das im öffentlichen Stromnetz integriert ist (AmpaCity in Essen). Dieses Semester ist ebenfalls eine Exkursion zum Beschleunigerzentrum DESY (Deutsches Elektronen-Synchrotron) geplant.

Eine weitere Veranstaltung in Münster ist das sogenannte "Meet your Prof". Dabei können Studenten des ersten Semesters (also ihr!) ihre Professoren näher kennen lernen. In gemütlicher Atmosphäre -- mit Orangensaft und Salami-Brötchen -- sitzen die Studenten mit ihren Profs am Tisch. Diese erzählen von ihrer aktuellen Forschung, aber auch Fragen zum Privatleben sind möglich und sogar erwünscht. Das ist ein Teil des Dialogs, den die jDPG möchte: Studenten treffen Professoren, Professoren treffen Studenten.

Von den Deutschlandweiten Veranstaltungen ist die Sommerexkursion das Highlight für die Mitglieder der jDPG. 5~Tage lang treffen sich Jungphysiker aus ganz Deutschland in einer Stadt, um Einblicke in die aktuelle Forschung und Entwicklung zu bekommen, und zwar nicht im Hörsaal, sondern direkt am Schauplatz des Geschehens. Neben dem fachlichen Teil ist aber auch eine Kneipentour und ein gemeinsames Grillen ein fester Bestandteil des Programmes.\\
Dieses Jahr ging es, passend zum internationalen Jahr des Lichts, nach Jena, die Stadt der Optik.

Besonders wichtig für die Kommunikation und Organisation der jDPG ist die jährliche Mitgliederversammlung. Dort werden wichtige Beschlüsse gefasst und der gesamte jDPG~Vorstand wird neugewählt.\\
Außerdem informieren der Newsletter auf der Internetseite und das Physik~Journal, die Mitgliederzeitschrift der Deutschen~Physikalischen~Gesellschaft, zusätzlich.

Zirka 60~aktive Mitglieder organisieren die Projekte der jDPG. Anna Bakenecker sieht in der Werbung neuer aktiver Mitglieder einer der wichtigsten Aufgaben der kommenden Jahre. Das Tolle an der jDPG sei, dass man in kurzer Zeit viel erreichen könne, sagt die Physikstudentin.
\begin{quote}
\textit{"In vielen Organisationen versickern neue Ideen. Bei uns werden sie umgesetzt."}
\end{quote}
Wer selbst aktiv werden möchte, kann dies in den zahlreichen Regionalgruppen vor Ort, oder auf bundesweiter Ebene in Arbeitskreisen.\\
Die Regionalgruppe Münster trifft sich einmal im Monat zum Stammtisch, zu dem auch du gerne vorbeischauen kannst, natürlich unverbindlich ;)

Der Anstoß zur Gründung der jDPG kam im Jahr~2005 aus der Deutschen~Physikalischen~Gesellschaft selbst. Diese hatte zwar 30~Prozent studentische Mitglieder, aber die Angebote für diese junge Zielgruppe fehlten. Als die DPG auf die Studenten der TU~Dresden zuging, haben sich spontan fünf Jungphysiker gefunden, die bereit waren, diese Lücke zu schließen. "Wir haben das damals für ein kleines Versuchsprojekt gehalten", erzählt jDPG-Mitglied René Pfitzner. Mit der rasanten Entwicklung der vergangenen Jahre habe damals niemand gerechnet.

Die erste Veranstaltung der jDPG war eine Sommerexkursion nach Berlin. Das Max-Born-Institut, das Albert-Einstein-Institut und auch die Firma Rolls~Royce standen auf dem Programm.

\begin{quote}
\textit{"Das hat uns einen richtigen Schub gegeben"},
\end{quote}
erinnert sich René Pfitzner. Aus ganz Deutschland hätten sich Freiwillige gemeldet, die die Idee der jDPG unterstützen wollten. So haben sich langsam die ersten Regionalgruppen gebildet und Ende~2006 wurde die jDPG offiziell gegründet.

Doch die steigenden Mitgliederzahlen zeigten auch, dass eine lose Arbeitsgruppe, wie sie am Anfang geplant war, nicht mehr ausreichen würde. Eine neue Struktur musste her und diese wurde auf der ersten Mitgliederversammlung im Jahr~2007 beschlossen und seitdem immer wieder verändert.

\begin{quote}
\textit{"Wir haben die jDPG nie auf dem Reißbrett geplant"},
\end{quote}
sagt der langjährige Bundessprecher Heinrich.
\begin{quote}
\textit{"Die Struktur ist immer aus den neuen Anforderungen gewachsen."}
\end{quote}

Heute ist die jDPG mit insgesamt 31~Regionalgruppen eine bundesweit aktive Organisation mit festem Programm.

Weitere Informationen findest du auf

\begin{center}
\textbf{\url{http://www.jdpg.de}}.
\end{center}

Die Angebote der Regionalgruppe~Münster und den Termin für den nächsten Stammtisch findest du auch direkt unter

\begin{center}
\textbf{\url{http://muenster.jdpg.de}}.
\end{center}

Bei weiteren Fragen kannst du dich auch gerne an René wenden

\begin{center}
\textbf{\email{henke.rene@wwu.de}}
\end{center}

\includegraphicscompressed[width=\columnwidth]{res/jdpg_foto.png}

\fibelsig{David Steffen, René Henke und Miriam Neumann}
\end{multicols}

\begin{center}
\includegraphics[width=0.8\textwidth]{res/comics/xkcd_experiment.png}
\end{center}